\section{Digital Signatures}

Providing privacy is independent from providing integrity. 
Digital signature schemes are public-key equivalent of MACs 
in private-key cryptography. A signer uses their private 
key $sk$ to “sign” a message, and anyone with the singer's 
public key can verify the integrity of that signature. 

Message authentication codes (MACs) require a shared secret key between parties.
Digital signatures avoid the need to establish a secret key with every single receiver.
The sender only computes a single signature that all recipients can verify.
Digital signatures are publicly verifiable: everyone can check they got the same
authenticated message.
Digital signatures are also transferable: signatures can be shown to a third party and
verified (important for certificates and public key infrastructure as we will see
later).
They also provide an important property of non-repudiation:
if a sender signs a message, they cannot later deny doing that, which is not true for MACs. 

A signature scheme consists of three PPT algorithms: 
\texttt{Gen}, \texttt{Sign}, \texttt{Verify} such that:
Key generation algorithm \texttt{Gen} takes as input a security 
parameter $1^n$ and outputs a pair of keys $(pk, sk)$.
Signing algorithm \texttt{Sign} takes as input a private key 
$sk$ and a message $m$ and
outputs a signature $\sigma = [m^d \mod N]$.
Verification algorithm \texttt{Verify} takes as input a 
public key $pk$, a message $m$, and a
signature $\sigma$. It outputs a bit $b$ either a 1 meaning valid or 0 meaning invalid:
$b = \text{\texttt{Verify}}_{pk}(m, \sigma)$.

An adversary should not be able to output a forgery 
even after seeing message/signature pairs. 

Similar to encryption, signature schemes are less efficient than MACs.
If we need to sign a long message, we can use a 
hybrid hash-and-sign construction where we generate a hash digest of 
the message and then sign the hash. 

The RSA signature scheme can be constructed as follows:
\begin{itemize}
    \item \texttt{Gen}: on input $1^n$ run \texttt{GenRSA} 
    to obtain $(N, e, d)$.
    \item \texttt{Sign}: on input a private key $sk = ⟨N, d⟩$ 
    and a message $m \in \mathbb{Z}_N^*$, compute signature 
    $\sigma = [m^d\ mod N]$.
    \item \texttt{Verify}: on input a public key $pk = ⟨N, e⟩$, 
    a message $m$, and a signature $\sigma \in \mathbb{Z}_N^*$, 
    output 1 iff $m = [\sigma^e \mod N]$.
\end{itemize}

An identification scheme is an interactive protocol that 
allows one party to prove
its identity to another.
We call the parties prover and verifier.
Verifier has the public key, prover convinces the verifier that it holds the
corresponding secret key.
